\chapter{Vladimir Voevodsky --- Fields Medal 2002}\label{ch2002vladimirvoevodsky:vladimir-voevodsky-fields-medal-2002-a-comprehensive-report-on-his-work-and-impact}

\begin{center}\rule{0.5\linewidth}{0.5pt}\end{center}

\textbf{Prepared:} July 2026 \textbf{Subject:} Vladimir Voevodsky (June
4, 1966 -- September 30, 2017, Moscow, Russia) \textbf{Position:}
Professor, Institute for Advanced Study, Princeton (2002--2017)
\textbf{Doctoral Advisor:} George Shabat (Moscow State University);
David Kazhdan (Harvard University) \textbf{Fields Medal Citation:}
\emph{``For his development of the homotopy theory of schemes and the
proof of the Milnor conjecture.''}

\begin{center}\rule{0.5\linewidth}{0.5pt}\end{center}

\section{1. Introduction}\label{ch2002vladimirvoevodsky:introduction}

Vladimir Voevodsky was awarded the Fields Medal in 2002 for his creation
of the \textbf{homotopy theory of schemes} and his proof of the
\textbf{Milnor conjecture}. These achievements transformed algebraic
geometry and number theory, providing a new framework for understanding
the relationship between algebraic cycles, cohomology theories, and the
arithmetic of fields.

\begin{prereqs} Algebraic geometry (motives, cohomology), homotopy theory (model categories, spectra), basic category theory. \end{prereqs}

\begin{objectives} Explain the univalence axiom and its implications; describe motivic cohomology; understand the proof of the Milnor conjecture. \end{objectives}


Voevodsky's work is characterized by notable originality. The
homotopy theory of schemes (A¹-homotopy theory) is a completely new
foundation for algebraic geometry, modeled on the homotopy theory of
topological spaces. Motivic cohomology, which he developed with Suslin,
Friedlander, and others, provides a universal cohomology theory for
algebraic varieties.

The proof of the Milnor conjecture, which relates the Milnor K-theory of
a field to its Galois cohomology, was a landmark in arithmetic geometry.
It was the first of many results in what became the \textbf{Bloch--Kato
conjecture} (now the \textbf{norm residue theorem}).

\begin{center}\rule{0.5\linewidth}{0.5pt}\end{center}

\section{2. Biographical Background}\label{ch2002vladimirvoevodsky:biographical-background}

\subsection{2.1 Early Life}\label{ch2002vladimirvoevodsky:early-life}

Vladimir Voevodsky was born in Moscow in 1966. His father was a
physicist and his mother a chemist. He showed mathematical talent from
an early age and was selected for the Kolmogorov boarding school at
Moscow State University.

\subsection{2.2 Education}\label{ch2002vladimirvoevodsky:education}

Voevodsky entered Moscow State University in 1983 but left before
completing his degree, finding the environment too restrictive. He
worked independently and later studied at Harvard University, where he
completed his PhD in 1992 under David Kazhdan. His thesis, ``Homology of
schemes and the Milnor conjecture,'' already contained the seeds of his
later work.

\subsection{2.3 Academic Career}\label{ch2002vladimirvoevodsky:academic-career}

\begin{itemize}
\tightlist
\item
 \textbf{1992--1995:} Postdoctoral fellow, Institute for Advanced
 Study, Princeton.
\item
 \textbf{1995--2002:} Associate Professor then Professor, Northwestern
 University.
\item
 \textbf{2002--2017:} Professor, Institute for Advanced Study,
 Princeton.
\item
 \textbf{2002:} Fields Medal, International Congress of Mathematicians,
 Beijing.
\item
 \textbf{2017:} Died from an aortic aneurysm at age 51.
\end{itemize}

\begin{center}\rule{0.5\linewidth}{0.5pt}\end{center}

\section{3. Motivic Cohomology}\label{ch2002vladimirvoevodsky:motivic-cohomology}

\subsection{3.1 The Problem}\label{ch2002vladimirvoevodsky:the-problem}

Before Voevodsky, algebraic geometry lacked a satisfactory theory of
\textbf{motivic cohomology} --- a universal cohomology theory for
algebraic varieties that would unify the various Weil cohomology
theories (singular cohomology, étale cohomology, de Rham cohomology,
etc.).

Grothendieck had conjectured the existence of such a theory (the theory
of \textbf{motives}), but constructing it remained an open problem.

\subsection{3.2 Voevodsky's Construction}\label{ch2002vladimirvoevodsky:voevodskys-construction}

Voevodsky, with Andrei Suslin and Eric Friedlander, constructed a theory
of motivic cohomology using: - \textbf{The category of presheaves with
transfers:} A category of presheaves that respect algebraic
correspondences. - \textbf{The Nisnevich topology:} A Grothendieck
topology that is finer than the Zariski topology but coarser than the
étale topology. - \textbf{A¹-homotopy invariance:} The property that the
cohomology is invariant under homotopy with respect to the affine line.

The resulting theory satisfies all the expected properties: - Long exact
sequences. - Mayer--Vietoris sequences. - Product structure. -
Comparison with K-theory.

\subsection{3.3 The Comparison with
K-Theory}\label{ch2002vladimirvoevodsky:the-comparison-with-k-theory}

The fundamental relationship between motivic cohomology and algebraic
K-theory is given by the \textbf{Atiyah--Hirzebruch spectral sequence}:

\[
E_2^{p,q} = H^{p-q}(X, \mathbb{Z}(-q)) \Rightarrow K_{-p-q}(X)
\]

This spectral sequence, proved by Voevodsky and Friedlander, relates
motivic cohomology to K-theory, providing an effective tool for computing
both.

\begin{center}\rule{0.5\linewidth}{0.5pt}\end{center}

\section{4. A¹-Homotopy Theory}\label{ch2002vladimirvoevodsky:auxb9-homotopy-theory}

\subsection{4.1 Motivation}\label{ch2002vladimirvoevodsky:motivation}

In topology, the homotopy category of spaces is obtained from the
category of topological spaces by inverting the interval \([0,1]\). In
algebraic geometry, Voevodsky's insight was to replace the interval by
the affine line \(\mathbb{A}^1\).

\subsection{4.2 The Construction}\label{ch2002vladimirvoevodsky:the-construction}

\textbf{A¹-homotopy theory} is the homotopy theory obtained from the
category of schemes by: 1. \textbf{Inverting the affine line:} Making
\(\mathbb{A}^1\) contractible. 2. \textbf{Using the Nisnevich topology:}
To handle the etale topology of schemes. 3. \textbf{Applying model
category techniques:} To obtain a well-behaved homotopy category.

The resulting category, called the \textbf{A¹-homotopy category} or the
\textbf{motivic homotopy category}, is a triangulated category that
contains: - The usual homotopy category of algebraic varieties. -
Motivic complexes. - The stable homotopy category of algebraic
varieties.

\subsection{4.3 The Motivic Sphere
Spectrum}\label{ch2002vladimirvoevodsky:the-motivic-sphere-spectrum}

The \textbf{motivic sphere spectrum} \(\mathbb{S}\) is the unit object
of the stable motivic homotopy category. Its homotopy groups are the
\textbf{motivic homotopy groups of spheres}, which are among the most
important objects in motivic homotopy theory.

\begin{center}\rule{0.5\linewidth}{0.5pt}\end{center}

\section{5. The Milnor Conjecture}\label{ch2002vladimirvoevodsky:the-milnor-conjecture}

\subsection{5.1 Background}\label{ch2002vladimirvoevodsky:background}

The \textbf{Milnor K-theory} \(K^M_*(F)\) of a field \(F\) is a graded
ring generated by symbols \(\{a_1, \ldots, a_n\}\) with
\(a_i \in F^\times\), subject to the Steinberg relation:

\[
\{a, 1-a\} = 0
\]

The \textbf{Galois cohomology} \(H^*(F, \mathbb{Z}/2\mathbb{Z})\) is the
cohomology of the absolute Galois group of \(F\) with coefficients in
\(\mathbb{Z}/2\mathbb{Z}\).

The \textbf{Milnor conjecture} (1970) states that there is an
isomorphism:

\[
K^M_n(F)/2 \cong H^n(F, \mathbb{Z}/2\mathbb{Z})
\]

\subsection{5.2 Voevodsky's Proof}\label{ch2002vladimirvoevodsky:voevodskys-proof}

Voevodsky proved the Milnor conjecture in 1996 using: 1. \textbf{Motivic
cohomology:} To relate Milnor K-theory to motivic cohomology. 2.
\textbf{The Steenrod algebra:} In motivic cohomology, to detect the
relevant classes. 3. \textbf{The norm residue homomorphism:} Constructed
by Tate and refined by others.

The proof is a tour de force that combines algebraic K-theory, motivic
cohomology, and the arithmetic of fields.

\subsection{5.3 The Map}\label{ch2002vladimirvoevodsky:the-map}

The isomorphism is given by the \textbf{norm residue map}:

\[
K^M_n(F)/2 \to H^n(F, \mathbb{Z}/2\mathbb{Z})
\]

sending a symbol \(\{a_1, \ldots, a_n\}\) to the cup product
\((a_1) \cup \cdots \cup (a_n)\), where \((a_i)\) is the class of
\(a_i\) in \(H^1(F, \mathbb{Z}/2\mathbb{Z}) = F^\times/(F^\times)^2\).

\begin{center}\rule{0.5\linewidth}{0.5pt}\end{center}

\section{6. The Bloch--Kato Conjecture}\label{ch2002vladimirvoevodsky:the-blochkato-conjecture}

\subsection{6.1 Statement}\label{ch2002vladimirvoevodsky:statement}

The \textbf{Bloch--Kato conjecture} (now the \textbf{norm residue
theorem}) generalizes the Milnor conjecture from
\(\mathbb{Z}/2\mathbb{Z}\) coefficients to \(\mathbb{Z}/l\mathbb{Z}\)
coefficients for any prime \(l\):

\[
K^M_n(F)/l \cong H^n(F, \mu_l^{\otimes n})
\]

where \(\mu_l\) is the group of \(l\)-th roots of unity.

\subsection{6.2 Voevodsky's Proof}\label{ch2002vladimirvoevodsky:voevodskys-proof-1}

Voevodsky, together with Markus Rost, proved the Bloch--Kato conjecture
in full generality. The proof, completed in 2011, runs to several
hundred pages and uses: - \textbf{Motivic cohomology and the Steenrod
algebra.} - \textbf{The theory of symmetric powers.} - \textbf{The norm
residue homomorphism.} - \textbf{The Rost--Voevodsky theorem on the
motivic cohomology of fields.}

\subsection{6.3 Impact}\label{ch2002vladimirvoevodsky:impact}

The norm residue theorem is one of the important results in
arithmetic geometry. It has implications for: - \textbf{Class field
theory:} The proof of the Bloch--Kato conjecture completes the program
of ``higher class field theory.'' - \textbf{Algebraic K-theory:} The
norm residue theorem is essential for computing K-groups of fields. -
\textbf{The cohomology of varieties:} The theorem provides a bridge
between K-theory and étale cohomology.

\begin{center}\rule{0.5\linewidth}{0.5pt}\end{center}

\section{7. The Steenrod Algebra in Motivic
Cohomology}\label{ch2002vladimirvoevodsky:the-steenrod-algebra-in-motivic-cohomology}

\subsection{7.1 The Classical Steenrod
Algebra}\label{ch2002vladimirvoevodsky:the-classical-steenrod-algebra}

In topology, the \textbf{Steenrod algebra} is the algebra of natural
operations on singular cohomology with \(\mathbb{F}_2\) coefficients. It
is generated by the Steenrod squares:

\[
\text{Sq}^i: H^n(X) \to H^{n+i}(X)
\]

\subsection{7.2 The Motivic Steenrod
Algebra}\label{ch2002vladimirvoevodsky:the-motivic-steenrod-algebra}

Voevodsky constructed an analog of the Steenrod algebra in motivic
cohomology. The motivic Steenrod operations act on motivic cohomology
and satisfy analogous properties: - \textbf{Cartan formula:}
\(\text{Sq}(x \cup y) = \text{Sq}(x) \cup \text{Sq}(y)\). - \textbf{Adem
relations:}
\(\text{Sq}^a \text{Sq}^b = \sum \binom{b-1}{a-2i} \text{Sq}^{a+b-i} \text{Sq}^i\).
- \textbf{Stability:} \(\text{Sq}^0 = 1\).

\subsection{7.3 Applications}\label{ch2002vladimirvoevodsky:applications}

The motivic Steenrod algebra is essential for: - \textbf{The Milnor
conjecture:} The detection theorem for motivic cohomology classes. -
\textbf{The Bloch--Kato conjecture:} The same techniques applied to
general primes. - \textbf{Computations in motivic cohomology:}
Determining the structure of the cohomology of fields.

\begin{center}\rule{0.5\linewidth}{0.5pt}\end{center}

\section{8. Transfers in Algebraic
K-Theory}\label{ch2002vladimirvoevodsky:transfers-in-algebraic-k-theory}

\subsection{8.1 The Problem}\label{ch2002vladimirvoevodsky:the-problem-1}

A \textbf{transfer} (or \textbf{norm}) in algebraic K-theory is a
homomorphism:

\[
\text{N}_{E/F}: K_*(E) \to K_*(F)
\]

associated to a finite field extension \(E/F\). The existence of
transfers is essential for many applications of K-theory.

\subsection{8.2 Voevodsky's
Construction}\label{ch2002vladimirvoevodsky:voevodskys-construction-1}

Voevodsky constructed transfers in algebraic K-theory for finite
morphisms of schemes. His construction uses: - \textbf{Algebraic
correspondences:} To define the transfer on the level of cycles. -
\textbf{The category of presheaves with transfers:} To extend the
transfer to general K-theory. - \textbf{The homotopy invariance of
K-theory:} To show that the transfer is well-defined.

\begin{center}\rule{0.5\linewidth}{0.5pt}\end{center}

\section{9. The Category of Presheaves with
Transfers}\label{ch2002vladimirvoevodsky:the-category-of-presheaves-with-transfers}

\subsection{9.1 Definition}\label{ch2002vladimirvoevodsky:definition}

A \textbf{presheaf with transfers} is a contravariant functor from the
category of finite correspondences (smooth schemes with algebraic
correspondences as morphisms) to abelian groups.

\subsection{9.2 Properties}\label{ch2002vladimirvoevodsky:properties}

\begin{itemize}
\tightlist
\item
 \textbf{Cohomological properties:} Sheaf cohomology in the Nisnevich
 topology gives motivic cohomology.
\item
 \textbf{A¹-invariance:} The presheaves can be A¹-localized to obtain
 motivic complexes.
\item
 \textbf{Tensor structure:} The category of presheaves with transfers
 is a tensor category.
\end{itemize}

\subsection{9.3 Role in Voevodsky's Work}\label{ch2002vladimirvoevodsky:role-in-voevodskys-work}

The category of presheaves with transfers is the foundation for
Voevodsky's construction of motivic cohomology. It provides the right
framework for: - Defining motivic complexes. - Proving the comparison
with K-theory. - Constructing the Atiyah--Hirzebruch spectral sequence.

\begin{center}\rule{0.5\linewidth}{0.5pt}\end{center}

\section{10. The Motivic Complexes}\label{ch2002vladimirvoevodsky:the-motivic-complexes}

\subsection{10.1 Definition}\label{ch2002vladimirvoevodsky:definition-1}

The \textbf{motivic complex} \(\mathbb{Z}(n)\) is a complex of
presheaves with transfers whose cohomology gives motivic cohomology:

\[
H^{p,q}(X, \mathbb{Z}) = \mathbb{H}^p(X, \mathbb{Z}(q))
\]

\subsection{10.2 Construction}\label{ch2002vladimirvoevodsky:construction}

Voevodsky constructs \(\mathbb{Z}(n)\) as: \[
\mathbb{Z}(n) = \mathbb{Z}(n-1) \otimes \mathbb{G}_m[-1]
\] where \(\mathbb{G}_m\) is the multiplicative group sheaf. This gives
a canonical construction of motivic complexes from the simplest case
\(n = 1\).

\subsection{10.3 Properties}\label{ch2002vladimirvoevodsky:properties-1}

\begin{itemize}
\tightlist
\item
 \(H^{p,q}(X, \mathbb{Z}) = 0\) for \(p > q + \dim(X)\).
\item
 \(H^{p,p}(\text{Spec }F, \mathbb{Z}) = K^M_p(F)\) for a field \(F\).
\item
 \(H^{p,q}(\text{Spec }F, \mathbb{Z}/l) = H^p(F, \mu_l^{\otimes q})\)
 for \(p \le q\).
\end{itemize}

\begin{center}\rule{0.5\linewidth}{0.5pt}\end{center}

\section{11. The Wild Kernel}\label{ch2002vladimirvoevodsky:the-wild-kernel}

\subsection{11.1 Definition}\label{ch2002vladimirvoevodsky:definition-2}

The \textbf{wild kernel} is the subgroup of the Bloch--Kato cohomology
that captures the difference between the motivic cohomology and the
étale cohomology of a field.

\subsection{11.2 Voevodsky's Work}\label{ch2002vladimirvoevodsky:voevodskys-work}

Voevodsky studied the wild kernel and showed that it vanishes for fields
of characteristic zero, completing the proof of the Bloch--Kato
conjecture for such fields.

\begin{center}\rule{0.5\linewidth}{0.5pt}\end{center}

\section{12. The Homotopy Fixed Point
Theorem}\label{ch2002vladimirvoevodsky:the-homotopy-fixed-point-theorem}

\subsection{12.1 Statement}\label{ch2002vladimirvoevodsky:statement-1}

The \textbf{homotopy fixed point theorem} for motivic cohomology states:

\[
H^{p,q}(X, \mathbb{Z}/l) \cong H^{p,q}(X_{ét}, \mu_l^{\otimes q})
\]

for smooth varieties over fields of characteristic not dividing \(l\).

\subsection{12.2 Proof}\label{ch2002vladimirvoevodsky:proof}

The proof uses: - The rigidity theorem for motivic cohomology. - The
comparison between the Nisnevich and étale topologies. - The theory of
Steenrod operations.

\begin{center}\rule{0.5\linewidth}{0.5pt}\end{center}

\section{13. The Rigidity Theorem}\label{ch2002vladimirvoevodsky:the-rigidity-theorem}

\subsection{13.1 Statement}\label{ch2002vladimirvoevodsky:statement-2}

The \textbf{rigidity theorem} states that for a Henselian local ring
\(R\) with residue field \(k\), the map:

\[
H^{p,q}(\text{Spec }R, \mathbb{Z}/l) \to H^{p,q}(\text{Spec }k, \mathbb{Z}/l)
\]

is an isomorphism for \(p \le q\).

\subsection{13.2 Applications}\label{ch2002vladimirvoevodsky:applications-1}

The rigidity theorem is essential for: - Computing motivic cohomology of
local fields. - Proving the homotopy fixed point theorem. - Relating
motivic cohomology to étale cohomology.

\begin{center}\rule{0.5\linewidth}{0.5pt}\end{center}

\section{14. The Comparison with Étale
Cohomology}\label{ch2002vladimirvoevodsky:the-comparison-with-uxe9tale-cohomology}

\subsection{14.1 The Comparison Map}\label{ch2002vladimirvoevodsky:the-comparison-map}

Voevodsky constructed a comparison map:

\[
H^{p,q}(X, \mathbb{Z}/l) \to H^p(X_{ét}, \mu_l^{\otimes q})
\]

\subsection{14.2 The Bloch--Kato Theorem}\label{ch2002vladimirvoevodsky:the-blochkato-theorem}

The Bloch--Kato conjecture (now theorem) implies that the comparison map
is an isomorphism for \(p \le q\). This is the \textbf{norm residue
isomorphism}.

\begin{center}\rule{0.5\linewidth}{0.5pt}\end{center}

\section{15. The Motivic Homotopy
Category}\label{ch2002vladimirvoevodsky:the-motivic-homotopy-category}

\subsection{15.1 The Stable Category}\label{ch2002vladimirvoevodsky:the-stable-category}

The \textbf{stable motivic homotopy category} \(\mathcal{SH}(S)\) is
obtained from the category of smooth schemes over \(S\) by: 1. Inverting
\(\mathbb{A}^1\). 2. Using the Nisnevich topology. 3. Stabilizing with
respect to the suspension by \(\mathbb{P}^1\).

\subsection{15.2 Properties}\label{ch2002vladimirvoevodsky:properties-2}

\begin{itemize}
\tightlist
\item
 \(\mathcal{SH}(S)\) is a triangulated tensor category.
\item
 It contains all known cohomology theories for algebraic varieties.
\item
 The unit object is the motivic sphere spectrum.
\item
 The endomorphism ring of the unit is the Grothendieck--Witt ring.
\end{itemize}

\subsection{15.3 Impact}\label{ch2002vladimirvoevodsky:impact-1}

The motivic homotopy category has become a central object in modern
algebraic geometry. It provides a framework for: - Studying the
cohomology of algebraic varieties. - Applying homotopy-theoretic methods
to algebraic geometry. - Understanding the relationship between
different cohomology theories.

\begin{center}\rule{0.5\linewidth}{0.5pt}\end{center}

\section{16. The Cross Product and the External
Product}\label{ch2002vladimirvoevodsky:the-cross-product-and-the-external-product}

\subsection{16.1 The Cross Product}\label{ch2002vladimirvoevodsky:the-cross-product}

The \textbf{cross product} in motivic cohomology:

\[
H^{p,q}(X, \mathbb{Z}) \otimes H^{p',q'}(Y, \mathbb{Z}) \to H^{p+p', q+q'}(X \times Y, \mathbb{Z})
\]

is defined using the external product of motivic complexes.

\subsection{16.2 The Cup Product}\label{ch2002vladimirvoevodsky:the-cup-product}

The cup product is obtained from the cross product by pulling back along
the diagonal:

\[
x \cup y = \Delta^*(x \times y)
\]

\begin{center}\rule{0.5\linewidth}{0.5pt}\end{center}

\section{17. The Motivic Cohomology of a
Field}\label{ch2002vladimirvoevodsky:the-motivic-cohomology-of-a-field}

\subsection{17.1 The Structure}\label{ch2002vladimirvoevodsky:the-structure}

For a field \(F\), the motivic cohomology ring is:

\[
H^{*,*}(\text{Spec }F, \mathbb{Z}) = K^M_*(F)
\]

the Milnor K-theory of \(F\). This is the key result that connects
motivic cohomology to K-theory.

\subsection{17.2 Voevodsky's Theorem}\label{ch2002vladimirvoevodsky:voevodskys-theorem}

Voevodsky proved that for a field \(F\), the canonical map:

\[
K^M_n(F) \to H^{n,n}(\text{Spec }F, \mathbb{Z})
\]

is an isomorphism. This is the \textbf{Suslin--Voevodsky theorem}.

\begin{center}\rule{0.5\linewidth}{0.5pt}\end{center}

\section{18. The Friedlander--Suslin--Voevodsky Theory of Motivic
Cohomology}\label{ch2002vladimirvoevodsky:the-friedlandersuslinvoevodsky-theory-of-motivic-cohomology}

\subsection{18.1 The Textbook}\label{ch2002vladimirvoevodsky:the-textbook}

Eric Friedlander, Andrei Suslin, and Vladimir Voevodsky co-authored a
foundational text, ``Cycles, Transfers, and Motivic Homology Theories,''
which provides a comprehensive treatment of motivic cohomology.

\subsection{18.2 The Standard Reference}\label{ch2002vladimirvoevodsky:the-standard-reference}

This text is the standard reference for: - The category of presheaves
with transfers. - Motivic complexes and their properties. - The
relationship between motivic cohomology and K-theory.

\begin{center}\rule{0.5\linewidth}{0.5pt}\end{center}

\section{19. The Proof of the Bloch--Kato
Conjecture}\label{ch2002vladimirvoevodsky:the-proof-of-the-blochkato-conjecture}

\subsection{19.1 The Strategy}\label{ch2002vladimirvoevodsky:the-strategy}

The proof of the Bloch--Kato conjecture proceeds in several steps: 1.
\textbf{Reduce to the prime \(l\) case.} The conjecture for general
\(l\) follows from the case where \(l\) is prime. 2. \textbf{Construct
the norm residue map.} Use Galois cohomology to construct the map from
Milnor K-theory to étale cohomology. 3. \textbf{Prove the map is
surjective.} Using the theory of Steenrod operations and motivic
cohomology. 4. \textbf{Prove the map is injective.} Using the rigidity
theorem and the comparison with étale cohomology.

\subsection{19.2 Rost's Contribution}\label{ch2002vladimirvoevodsky:rosts-contribution}

Markus Rost made essential contributions to the proof, particularly
concerning the \textbf{Rost cycle lemma} and the theory of \textbf{Rost
motives}.

\begin{center}\rule{0.5\linewidth}{0.5pt}\end{center}

\section{20. The Work of Voevodsky on Computer
Proofs}\label{ch2002vladimirvoevodsky:the-work-of-voevodsky-on-computer-proofs}

\subsection{20.1 The ``Mathematical
Repository''}\label{ch2002vladimirvoevodsky:the-mathematical-repository}

In his later years, Voevodsky became interested in computer verification
of mathematical proofs. He believed that the increasing complexity of
mathematics required formal verification to ensure correctness.

\subsection{20.2 The UniMath Project}\label{ch2002vladimirvoevodsky:the-unimath-project}

Voevodsky initiated the \textbf{UniMath} project, which aims to develop
a unified mathematical library in the Coq proof assistant. The project
is based on \textbf{homotopy type theory} (HoTT), a new foundation for
mathematics that Voevodsky helped develop.

\subsection{20.3 Homotopy Type Theory}\label{ch2002vladimirvoevodsky:homotopy-type-theory}

\textbf{Homotopy type theory} is a new foundation for mathematics based
on the analogy between types and homotopy types. Voevodsky made
important contributions to this theory, including the \textbf{univalence
axiom} and the \textbf{Voevodsky universe}.

\begin{center}\rule{0.5\linewidth}{0.5pt}\end{center}

\section{21. Impact on Mathematics}\label{ch2002vladimirvoevodsky:impact-on-mathematics}

\subsection{21.1 Motivic Cohomology}\label{ch2002vladimirvoevodsky:motivic-cohomology-1}

Voevodsky's construction of motivic cohomology is now a standard tool in
algebraic geometry. It is used for: - \textbf{Computing K-groups:} Using
spectral sequences. - \textbf{The arithmetic of varieties:} The
relationship between K-theory and étale cohomology. - \textbf{Birational
geometry:} The cohomology of algebraic varieties.

\subsection{21.2 The Bloch--Kato
Conjecture}\label{ch2002vladimirvoevodsky:the-blochkato-conjecture-1}

The proof of the Bloch--Kato conjecture is a landmark in arithmetic
geometry. It has implications for: - \textbf{Class field theory:} Higher
class field theory for fields. - \textbf{Algebraic K-theory:} The
computation of K-groups of fields. - \textbf{Galois cohomology:} The
structure of Galois cohomology groups.

\subsection{21.3 A¹-Homotopy Theory}\label{ch2002vladimirvoevodsky:auxb9-homotopy-theory-1}

A¹-homotopy theory has become an active area of research, with
applications to: - \textbf{The cohomology of algebraic varieties.} -
\textbf{Birational geometry.} - \textbf{Representation theory.} -
\textbf{Number theory.}

\subsection{21.4 Computer Verification}\label{ch2002vladimirvoevodsky:computer-verification}

Voevodsky's work on homotopy type theory and computer verification has
had a significant impact on the foundations of mathematics.

\begin{center}\rule{0.5\linewidth}{0.5pt}\end{center}

\section{22. Awards and Recognition}\label{ch2002vladimirvoevodsky:awards-and-recognition}

\textbf{2002:} - Fields Medal (International Mathematical Union)

\textbf{2001:} - Prize of the European Mathematical Society

\textbf{1997:} - Clay Research Award (with A. Suslin)

\textbf{1996:} - Prix de la Société Mathématique de France

\begin{center}\rule{0.5\linewidth}{0.5pt}\end{center}

\section{23. Selected Major
Publications}\label{ch2002vladimirvoevodsky:selected-major-publications}

\begin{enumerate}
\def\labelenumi{\arabic{enumi}.}
\item
 \textbf{``A¹-homotopy theory of schemes.''} \emph{Publications
 Mathématiques de l'IHÉS} 98, 2003, 1--48. --- \emph{The foundational
 paper on A¹-homotopy theory.}
\item
 \textbf{``Motivic cohomology groups are isomorphic to higher Chow
 groups''} (with E. Friedlander and A. Suslin). \emph{Annals of
 Mathematics} 148, 1998, 469--536. --- \emph{The comparison between
 motivic cohomology and Chow groups.}
\item
 \textbf{``The Milnor conjecture.''} \emph{Annals of Mathematics} 147,
 1998, 313--354. --- \emph{Proof of the Milnor conjecture on Galois
 cohomology.}
\item
 \textbf{``The Bloch--Kato conjecture.''} \emph{Journal of the American
 Mathematical Society} 24, 2011, 957--985. --- \emph{The proof of the
 Bloch--Kato conjecture (with M. Rost).}
\item
 \textbf{``Triangulated categories of motives over a field.''}
 \emph{Cycles, Transfers, and Motivic Homology Theories}, Princeton
 University Press, 2000. --- \emph{The construction of the triangulated
 category of motives.}
\item
 \textbf{``On the Steenrod operations in motivic cohomology.''}
 \emph{Izvestiya Mathematics} 65, 2001, 87--107. --- \emph{The
 construction of the motivic Steenrod algebra.}
\item
 \textbf{``The homotopy theory of simplicial sheaves.''}
 \emph{Compositio Mathematica} 120, 2000, 87--114. --- \emph{The
 foundations of A¹-homotopy theory.}
\end{enumerate}

\begin{center}\rule{0.5\linewidth}{0.5pt}\end{center}

\section{24. Selected Publications on Homotopy Type
Theory}\label{ch2002vladimirvoevodsky:selected-publications-on-homotopy-type-theory}

\begin{enumerate}
\def\labelenumi{\arabic{enumi}.}
\item
 \textbf{``The univalence axiom.''} \emph{Homotopy Type Theory:
 Univalent Foundations of Mathematics}, IAS, 2013. --- \emph{The
 univalence axiom and its consequences.}
\item
 \textbf{``A very short note on homotopy λ-calculus.''}
 \emph{Unpublished note}, 2006. --- \emph{The relationship between
 homotopy theory and type theory.}
\item
 \textbf{``The Coq proof assistant and homotopy type theory.''}
 \emph{Formal verification}, 2014.
\end{enumerate}

\begin{center}\rule{0.5\linewidth}{0.5pt}\end{center}

\section{25. Voevodsky's Mathematical
Style}\label{ch2002vladimirvoevodsky:voevodskys-mathematical-style}

\subsection{25.1 Originality}\label{ch2002vladimirvoevodsky:originality}

Voevodsky's work is characterized by notable originality. He did
not follow existing trends but created entirely new fields: -
A¹-homotopy theory. - Motivic cohomology. - The univalence axiom.

\subsection{25.2 Depth}\label{ch2002vladimirvoevodsky:depth}

His proofs are deep and systematic, often running to hundreds of pages.
He was not satisfied with partial results but insisted on complete
solutions.

\subsection{25.3 Vision}\label{ch2002vladimirvoevodsky:vision}

Voevodsky had a grand vision for mathematics: the unification of
algebraic geometry, topology, and number theory through motivic
cohomology, and the formal verification of mathematics through homotopy
type theory.

\begin{center}\rule{0.5\linewidth}{0.5pt}\end{center}

\section{26. The Legacy of Voevodsky}\label{ch2002vladimirvoevodsky:the-legacy-of-voevodsky}

\subsection{26.1 Motivic Cohomology}\label{ch2002vladimirvoevodsky:motivic-cohomology-2}

Voevodsky's construction of motivic cohomology is now a standard tool in
algebraic geometry. Every algebraic geometer working on the cohomology
of varieties uses his ideas.

\subsection{26.2 The Bloch--Kato Theorem}\label{ch2002vladimirvoevodsky:the-blochkato-theorem-1}

The proof of the Bloch--Kato conjecture is one of the great achievements
of arithmetic geometry. It solved a problem that had been open for over
30 years.

\subsection{26.3 Homotopy Type Theory}\label{ch2002vladimirvoevodsky:homotopy-type-theory-1}

Voevodsky's work on homotopy type theory may ultimately have the
greatest impact. It provides a new foundation for mathematics that is
compatible with computer verification.

\begin{center}\rule{0.5\linewidth}{0.5pt}\end{center}

\section{27. The Relationship with Other Cohomology
Theories}\label{ch2002vladimirvoevodsky:the-relationship-with-other-cohomology-theories}

\subsection{27.1 The Comparison with de Rham
Cohomology}\label{ch2002vladimirvoevodsky:the-comparison-with-de-rham-cohomology}

Motivic cohomology maps to de Rham cohomology through the
\textbf{Bloch--Ogus--Gabber theorem}. This comparison is essential for
understanding the relationship between algebraic cycles and differential
forms.

\subsection{27.2 The Comparison with Étale
Cohomology}\label{ch2002vladimirvoevodsky:the-comparison-with-uxe9tale-cohomology-1}

The map from motivic cohomology to étale cohomology is the \textbf{cycle
class map}. The Bloch--Kato theorem states that this map is an
isomorphism for torsion coefficients in the appropriate range.

\subsection{27.3 The Comparison with Singular
Cohomology}\label{ch2002vladimirvoevodsky:the-comparison-with-singular-cohomology}

For complex varieties, motivic cohomology maps to singular cohomology
via the \textbf{realization functor}. This functor sends the motivic
sphere to the usual sphere, and the A¹-homotopy category to the usual
homotopy category.

\begin{center}\rule{0.5\linewidth}{0.5pt}\end{center}

\section{28. The Motivic Cohomology of a
Point}\label{ch2002vladimirvoevodsky:the-motivic-cohomology-of-a-point}

\subsection{28.1 The Calculation}\label{ch2002vladimirvoevodsky:the-calculation}

For a field \(F\), the motivic cohomology of a point (Spec \(F\)) is:

\[
H^{p,q}(\text{Spec }F, \mathbb{Z}) =
\begin{cases}
K^M_q(F) & p = q \\
0 & p \neq q
\end{cases}
\]

\subsection{28.2 The Structure}\label{ch2002vladimirvoevodsky:the-structure-1}

The motivic cohomology ring of a field is therefore the Milnor K-theory
ring. This is the fundamental calculation that underlies the proof of
the Milnor conjecture.

\begin{center}\rule{0.5\linewidth}{0.5pt}\end{center}

\section{29. The Motivic Cohomology of
Curves}\label{ch2002vladimirvoevodsky:the-motivic-cohomology-of-curves}

\subsection{29.1 The Calculation}\label{ch2002vladimirvoevodsky:the-calculation-1}

For a smooth projective curve \(C\) over a field \(F\), the motivic
cohomology is:

\[
H^{p,q}(C, \mathbb{Z}) =
\begin{cases}
\mathbb{Z} & p = 0, q = 0 \\
F^\times & p = 1, q = 1 \\
\text{Pic}(C) & p = 2, q = 2 \\
K^M_0(F) & p = 0, q = 0
\end{cases}
\]

\subsection{29.2 Applications}\label{ch2002vladimirvoevodsky:applications-2}

This calculation is used in the study of algebraic cycles on curves and
in the theory of motives of curves.

\begin{center}\rule{0.5\linewidth}{0.5pt}\end{center}

\section{30. The Motivic Cohomology of Algebraic
Groups}\label{ch2002vladimirvoevodsky:the-motivic-cohomology-of-algebraic-groups}

\subsection{30.1 The General Theory}\label{ch2002vladimirvoevodsky:the-general-theory}

The motivic cohomology of algebraic groups can be computed using the
\textbf{classifying space} construction in the A¹-homotopy category.

\subsection{30.2 Applications}\label{ch2002vladimirvoevodsky:applications-3}

These calculations are essential for: - The theory of algebraic groups.
- The study of homogeneous varieties. - The theory of Chow groups and
algebraic cycles.

\begin{center}\rule{0.5\linewidth}{0.5pt}\end{center}

\section{31. The Motivic Cohomology of the Affine
Line}\label{ch2002vladimirvoevodsky:the-motivic-cohomology-of-the-affine-line}

\subsection{31.1 The Calculation}\label{ch2002vladimirvoevodsky:the-calculation-2}

The motivic cohomology of the affine line is:

\[
H^{p,q}(\mathbb{A}^1, \mathbb{Z}) = H^{p,q}(\text{Spec }F, \mathbb{Z})
\]

This is the \textbf{A¹-invariance} property of motivic cohomology.

\subsection{31.2 Consequences}\label{ch2002vladimirvoevodsky:consequences}

A¹-invariance implies that the motivic cohomology of any A¹-contractible
variety is the same as that of a point.

\begin{center}\rule{0.5\linewidth}{0.5pt}\end{center}

\section{32. Voevodsky's Work on Quadratic
Forms}\label{ch2002vladimirvoevodsky:voevodskys-work-on-quadratic-forms}

\subsection{32.1 The Connection}\label{ch2002vladimirvoevodsky:the-connection}

Voevodsky's proof of the Milnor conjecture has important implications
for the theory of quadratic forms. The Milnor conjecture relates the
Milnor K-theory of a field to its Galois cohomology, and this
relationship is used to study quadratic forms.

\subsection{32.2 The Milnor--Scharlau--Witt
Theorem}\label{ch2002vladimirvoevodsky:the-milnorscharlauwitt-theorem}

The Milnor conjecture implies the \textbf{Milnor--Scharlau--Witt
theorem}, which relates the Witt ring of quadratic forms to the étale
cohomology of the field.

\begin{center}\rule{0.5\linewidth}{0.5pt}\end{center}

\section{33. The Motivic Stable Homotopy
Category}\label{ch2002vladimirvoevodsky:the-motivic-stable-homotopy-category}

\subsection{33.1 The Construction}\label{ch2002vladimirvoevodsky:the-construction-1}

The \textbf{stable motivic homotopy category} is obtained from the
A¹-homotopy category by inverting the suspension by \(\mathbb{P}^1\).
This is the category of \textbf{motivic spectra}.

\subsection{33.2 Properties}\label{ch2002vladimirvoevodsky:properties-3}

\begin{itemize}
\tightlist
\item
 The category is symmetric monoidal.
\item
 It contains the motivic sphere spectrum.
\item
 Cohomology theories are represented by spectra.
\item
 The category is the universal target for orientable cohomology
 theories.
\end{itemize}

\begin{center}\rule{0.5\linewidth}{0.5pt}\end{center}

\section{34. The Motivic Thom Spectrum}\label{ch2002vladimirvoevodsky:the-motivic-thom-spectrum}

\subsection{34.1 Definition}\label{ch2002vladimirvoevodsky:definition-3}

The \textbf{motivic Thom spectrum} MGL is the algebraic analog of the
complex cobordism spectrum MU. It is the universal oriented cohomology
theory in algebraic geometry.

\subsection{34.2 Properties}\label{ch2002vladimirvoevodsky:properties-4}

\begin{itemize}
\tightlist
\item
 MGL has a theory of characteristic classes (Chern classes).
\item
 The MGL-cohomology of a point is the Lazard ring.
\item
 MGL is related to algebraic cobordism.
\end{itemize}

\begin{center}\rule{0.5\linewidth}{0.5pt}\end{center}

\section{35. Voevodsky's Impact on Algebraic
Topology}\label{ch2002vladimirvoevodsky:voevodskys-impact-on-algebraic-topology}

\subsection{35.1 The Motivic Adams Spectral
Sequence}\label{ch2002vladimirvoevodsky:the-motivic-adams-spectral-sequence}

The \textbf{motivic Adams spectral sequence} is a tool for computing the
motivic homotopy groups of spheres. It is the algebraic geometry analog
of the classical Adams spectral sequence.

\subsection{35.2 New Results on the Homotopy Groups of
Spheres}\label{ch2002vladimirvoevodsky:new-results-on-the-homotopy-groups-of-spheres}

The motivic Adams spectral sequence has produced new results on the
classical homotopy groups of spheres, using the motivic filtration to
simplify computations.

\begin{center}\rule{0.5\linewidth}{0.5pt}\end{center}

\section{36. The Univalence Axiom}\label{ch2002vladimirvoevodsky:the-univalence-axiom}

\subsection{36.1 Statement}\label{ch2002vladimirvoevodsky:statement-3}

The \textbf{univalence axiom} states that equivalence between types is
equivalent to equality between types:

\[
(A \simeq B) \simeq (A = B)
\]

\subsection{36.2 Consequences}\label{ch2002vladimirvoevodsky:consequences-1}

The univalence axiom is the key innovation of homotopy type theory. It
implies that isomorphic mathematical structures can be treated as
identical, which is essential for the practicality of HoTT.

\begin{center}\rule{0.5\linewidth}{0.5pt}\end{center}

\section{37. Conclusion}\label{ch2002vladimirvoevodsky:conclusion}

Vladimir Voevodsky was one of the most original mathematicians of his
generation. His creation of motivic cohomology and A¹-homotopy theory
transformed algebraic geometry. His proof of the Milnor conjecture and
the Bloch--Kato conjecture solved some of the important problems in
arithmetic geometry.

The Fields Medal citation --- ``for his development of the homotopy
theory of schemes and the proof of the Milnor conjecture'' --- captures
only part of his achievement. He went on to prove the Bloch--Kato
conjecture, develop homotopy type theory, and pioneer the formal
verification of mathematics.

His death in 2017 at age 51 cut short a career that had already changed
the face of mathematics. His ideas continue to influence algebraic
geometry, number theory, and the foundations of mathematics.

\begin{center}\rule{0.5\linewidth}{0.5pt}\end{center}

\begin{center}\rule{0.5\linewidth}{0.5pt}\end{center}

\section*{Exercises}
\addcontentsline{toc}{section}{Exercises}

\begin{exercise} Show that the Milnor K-theory of a field surjects onto its Galois cohomology with $\\mathbb{Z}/2$ coefficients. \end{exercise}

\begin{exercise} Verify the Bloch-Kato conjecture for $n = 1$: $K_n^M(F)/2 \\cong H^n_{\\text{\\'{e}t}}(F, \\mathbb{Z}/2)$. \end{exercise}

\begin{exercise} Construct the motivic stable homotopy category and verify it has the correct notion of weak equivalences. \end{exercise}

\begin{exercise} Prove that the univalence axiom implies functional extensionality in Homotopy Type Theory. \end{exercise}

\begin{center}\rule{0.5\linewidth}{0.5pt}\end{center}

\section*{Summary}
\addcontentsline{toc}{section}{Summary}

Voevodsky developed motivic cohomology and the theory of mixed motives, proving the Milnor conjecture and the Bloch-Kato conjecture. His work on Homotopy Type Theory and the univalence axiom created a new foundation for mathematics based on homotopy theory, with profound implications for computer-verified proofs.

\section{References and Further
Reading}\label{ch2002vladimirvoevodsky:references-and-further-reading}

\subsection{Official Sources}\label{ch2002vladimirvoevodsky:official-sources}

\begin{itemize}
\tightlist
\item
 IMU Fields Medal 2002 Announcement:
 https://www.mathunion.org/imu-awards/fields-medal/fields-medal-2002
\item
 Institute for Advanced Study: ``Vladimir Voevodsky, 1966--2017.''
\end{itemize}

\subsection{Key Technical References}\label{ch2002vladimirvoevodsky:key-technical-references}

\begin{itemize}
\tightlist
\item
 Voevodsky, V. ``A¹-homotopy theory of schemes.'' \emph{IHÉS
 Publications} 98, 2003.
\item
 Voevodsky, V. ``The Milnor conjecture.'' \emph{Annals of Mathematics}
 147, 1998.
\item
 Voevodsky, V. ``Triangulated categories of motives over a field.'' In
 \emph{Cycles, Transfers, and Motivic Homology Theories}, Princeton,
 2000.
\item
 Friedlander, E., Suslin, A., and Voevodsky, V. \emph{Cycles,
 Transfers, and Motivic Homology Theories}. Princeton University Press,
 2000.
\item
 Beilinson, A. ``Remarks on the Bloch--Kato conjecture.'' \emph{Journal
 of K-Theory} 7, 2011.
\end{itemize}

\subsection{Biography and Overview}\label{ch2002vladimirvoevodsky:biography-and-overview}

\begin{itemize}
\tightlist
\item
 Morel, F. ``The Work of Vladimir Voevodsky.'' ICM 2002 Laudatio.
\item
 Friedlander, E. ``Vladimir Voevodsky (1966--2017).'' \emph{Notices of
 the AMS}, 2018.
\item
 Deligne, P. ``Vladimir Voevodsky.'' \emph{IHÉS News}, 2017.
\item
 Mochizuki, S. ``On Voevodsky's proof of the Milnor conjecture.''
 \emph{K-Theory} 17, 1999.
\end{itemize}

\subsection{Textbooks}\label{ch2002vladimirvoevodsky:textbooks}

\begin{itemize}
\tightlist
\item
 Levine, M. \emph{Motivic Cohomology}. Clay Mathematics Monographs,
 2007.
\item
 Mazza, C., Voevodsky, V., and Weibel, C. \emph{Lectures on Motivic
 Cohomology}. Clay Mathematics Monographs, 2006.
\item
 Morel, F. \emph{A¹-Algebraic Topology over a Field}. Springer, 2012.
\item
 Asok, A. and Fasel, J. \emph{A¹-Homotopy Theory and the Unstable
 Classification of Vector Bundles}. Springer, 2014.
\end{itemize}

\begin{center}\rule{0.5\linewidth}{0.5pt}\end{center}


\begin{center}\rule{0.5\linewidth}{0.5pt}\end{center}

\section*{See Also}
\addcontentsline{toc}{section}{See Also}

This chapter relates to the following work by other Fields Medalists:
\begin{itemize}
  \item \textbf{\hyperref[ch1966alexandergrothendieck:alexander-grothendieck-fields-medal-1966-a-comprehensive-report]{See Chapter on related work}} --- Voevodsky extended Grothendieck's motivic cohomology.
  \item \textbf{\hyperref[ch2018peterscholze:peter-scholze-fields-medal-2018-a-comprehensive-report-on-his-work-and-impact]{See Chapter on related work}} --- Scholze's work on perfectoid spaces connects to Voevodsky.
  \item \textbf{\hyperref[ch1978danielg.quillen:daniel-gray-quillen-fields-medal-1978-a-comprehensive-report-on-his-work-and-impact]{See Chapter on related work}} --- Quillen's K-theory relates to motivic cohomology.
\end{itemize}

\emph{(End of report)}
